\documentclass[12pt,a4paper]{article}
\usepackage[utf8]{inputenc}
\usepackage{geometry}
\usepackage{xcolor}
\usepackage{listings}
\usepackage{graphicx}
\usepackage{hyperref}

\geometry{margin=1in}

\definecolor{codegreen}{rgb}{0,0.6,0}
\definecolor{codegray}{rgb}{0.5,0.5,0.5}
\definecolor{codepurple}{rgb}{0.58,0,0.82}
\definecolor{backcolour}{rgb}{0.95,0.95,0.92}

\lstdefinestyle{mystyle}{
    backgroundcolor=\color{backcolour},   
    commentstyle=\color{codegreen},
    keywordstyle=\color{magenta},
    numberstyle=\tiny\color{codegray},
    stringstyle=\color{codepurple},
    basicstyle=\ttfamily\footnotesize,
    breakatwhitespace=false,         
    breaklines=true,                 
    captionpos=b,                    
    keepspaces=true,                 
    numbers=left,                    
    numbersep=5pt,                  
    showspaces=false,                
    showstringspaces=false,
    showtabs=false,                  
    tabsize=2
}

\lstset{style=mystyle}

\title{Milestone 4 Report: Porting llama2.c to xv6}
\author{OS Fall 2025 Project Interstellar}
\date{\today}

\begin{document}

\maketitle

\section{Integration Strategy}

\subsection{Integration Approach}
The integration of \texttt{llama2.c} into xv6 involved porting the inference engine to run as a user-space program. Due to xv6's limited standard library and file system constraints, we adopted a hybrid approach:
\begin{enumerate}
    \item \textbf{Network-based Loading}: Instead of reading the 60MB model file from the xv6 file system (which is slow and size-constrained), we implemented a UDP client to fetch model weights and tokenizer data chunk-by-chunk from a host server.
    \item \textbf{Standard Library Emulation}: We replaced missing C standard library functions (e.g., \texttt{fseek}, \texttt{ftell}, \texttt{mmap}, \texttt{time}) with xv6 equivalents or custom implementations.
    \item \textbf{Hardware Acceleration}: We utilized RISC-V specific instructions (\texttt{rdtime}) for high-precision timing to measure token generation speed.
\end{enumerate}

\subsection{Modified Functions}
Key functions modified from the original \texttt{llama2.c}:
\begin{itemize}
    \item \textbf{\texttt{main}}: Completely rewritten to handle xv6 command-line arguments (which lack sophisticated parsing) and to initialize the network fetch process.
    \item \textbf{\texttt{read\_checkpoint}}: Refactored to use \texttt{fetch\_model\_weights} (UDP) instead of \texttt{fopen}/\texttt{fread}/\texttt{mmap}.
    \item \textbf{\texttt{build\_tokenizer}}: Refactored to use \texttt{fetch\_tokenizer} (UDP) instead of file I/O.
    \item \textbf{\texttt{time\_in\_ms}}: Modified to use the \texttt{rdcycle} (mapped to \texttt{rdtime}) system call for timing.
    \item \textbf{\texttt{error\_usage}}: Updated to reflect the simplified usage in xv6.
\end{itemize}

\subsection{Memory Management Changes}
The original \texttt{llama2.c} used \texttt{mmap} to map the model file into memory. xv6 does not support \texttt{mmap}.

\textbf{Change}: We replaced \texttt{mmap} with \texttt{malloc}.
\begin{enumerate}
    \item \textbf{Allocation}: We allocate a contiguous block of memory using \texttt{malloc} (backed by \texttt{sbrk} in xv6) large enough to hold the model parameters.
    \item \textbf{Loading}: We implemented \texttt{fetch\_model\_weights} which receives UDP packets and copies the data directly into this allocated memory buffer.
\end{enumerate}

\subsubsection{Code Snippet: Updated \texttt{read\_checkpoint} (Network Fetch)}
\begin{lstlisting}[language=C]
void read_checkpoint(char *checkpoint, Config *config, TransformerWeights *weights,
                     int *fd, float** data, ssize_t* file_size) {
    // 1. Fetch metadata and allocate memory
    // Note: fetch_model_weights handles the UDP communication
    // and returns a pointer to the fully loaded data in memory.
    *data = (float*) fetch_model_weights(0); // ID 0 for model
    if (*data == NULL) {
        fprintf(2, "Failed to fetch model weights\n");
        exit(1);
    }

    // 2. Parse header (Config) from the start of the buffer
    int *ptr = (int*)(*data);
    config->dim = ptr[0];
    config->hidden_dim = ptr[1];
    config->n_layers = ptr[2];
    config->n_heads = ptr[3];
    config->n_kv_heads = ptr[4];
    config->vocab_size = ptr[5];
    config->seq_len = ptr[6];

    // 3. Set up weight pointers into the buffer
    float* w = *data + 7; // Skip header
    weights->token_embedding_table = w;
    // ... (rest of weight assignments)
}
\end{lstlisting}

\subsubsection{Code Snippet: Updated \texttt{build\_tokenizer}}
\begin{lstlisting}[language=C]
void build_tokenizer(Tokenizer* t, char* tokenizer_path, int vocab_size) {
    // Fetch tokenizer data via UDP (ID 1)
    unsigned char* data = (unsigned char*) fetch_tokenizer(1);
    
    t->vocab_size = vocab_size;
    t->vocab = (char**)calloc(vocab_size, sizeof(char*));
    
    // Parse the fetched buffer
    int offset = 0;
    for (int i = 0; i < vocab_size; i++) {
        t->vocab[i] = (char*) (data + offset);
        int len = strlen(t->vocab[i]);
        offset += len + 1;
    }
}
\end{lstlisting}

\section{Verification \& Output}

\subsection{Determinism Test}
\textbf{Command}: \texttt{llm -s 42 -n 16 -i "Once upon a time"}

\textbf{Run 1 Output}:
\begin{verbatim}
Once upon a time, there was a little girl named Lily. She loved to play with her...
\end{verbatim}

\textbf{Run 2 Output}:
\begin{verbatim}
Once upon a time, there was a little girl named Lily. She loved to play with her...
\end{verbatim}

\textit{Observation}: The outputs are identical, confirming that the random number generator (seeded with 42) and the model inference are deterministic.

\subsection{Test Prompt Snippets}

\textbf{Prompt 1}: "Once upon a time"
\begin{verbatim}
Once upon a time, there was a little girl named Lily. She loved to play with her friends in the park. One day, she saw a big, red ball...
\end{verbatim}

\textbf{Prompt 2}: "The tiny dragon"
\begin{verbatim}
The tiny dragon was very sad. He wanted to fly, but his wings were too small. "I wish I could fly like the big dragons," he said...
\end{verbatim}

\textbf{Prompt 3}: "Hello world"
    \item \textbf{Symptom}: xv6 crashed with \texttt{panic: kerneltrap} and \texttt{scause=0x2} (Illegal Instruction) when running \texttt{llm}.
    \item \textbf{Root Cause}: The code used the \texttt{rdcycle} assembly instruction, which was not supported or enabled in our QEMU RISC-V configuration.
    \item \textbf{Solution}: We modified \texttt{kernel/riscv.h} and \texttt{kernel/sysproc.c} to use the \texttt{rdtime} pseudo-instruction instead, which reads the memory-mapped timer and is widely supported.
\end{itemize}

\subsection{Challenge 3: FPU Page Fault (\texttt{scause=0xf})}
\begin{itemize}
    \item \textbf{Symptom}: The \texttt{llm} program crashed with a "store page fault" (\texttt{scause=0xf}) pointing to FPU instructions.
    \item \textbf{Root Cause}: The RISC-V FPU is disabled by default. When the kernel returns to user mode, the \texttt{sstatus} register's FS (Floating Point Status) bits were set to "Off", causing any floating-point instruction to trap.
    \item \textbf{Solution}: We modified \texttt{kernel/trap.c} in the \texttt{prepare\_return} function (and \texttt{usertrap}) to explicitly enable the FPU by setting the FS bits to "Initial" (0x1) in \texttt{sstatus} before returning to user mode.
\end{itemize}

\subsection{Challenge 4: Missing Float Printing}
\begin{itemize}
    \item \textbf{Symptom}: The performance metric "tokens per second" was not printing correctly (garbage output).
    \item \textbf{Root Cause}: The minimal \texttt{printf} implementation in xv6 does not support the \texttt{\%f} format specifier.
    \item \textbf{Solution}: We implemented a helper function \texttt{print\_float} in \texttt{llm.c} that manually formats the floating-point number into integer and fractional parts for printing.
\end{itemize}

\section{Conclusion}

\subsection{Summary}
We successfully ported the \texttt{llama2.c} inference engine to the xv6 operating system. The project demonstrated that a modern, complex workload like an LLM can run on a simple educational OS by overcoming significant constraints in memory, I/O, and standard library support. Key achievements include implementing a UDP-based file loader, enabling hardware floating-point support in the kernel, and adapting the inference code to a bare-bones environment.

\subsection{Readiness for Milestone 5}
The system is fully ready for Milestone 5. The sequential implementation is verified and robust. The codebase is now set up for parallelization experiments:
\begin{itemize}
    \item \textbf{Profiling}: The \texttt{rdcycle} (via \texttt{rdtime}) system call provides the necessary timing infrastructure.
    \item \textbf{Workload}: The matrix multiplication loops in \texttt{llm.c} are identified as the primary targets for parallelization.
    \item \textbf{Stability}: The FPU and memory issues have been resolved, providing a stable baseline.
\end{itemize}

\subsection{Lessons Learned}
\begin{itemize}
    \item \textbf{Kernel-User Boundary}: Understanding how the kernel manages hardware state (like FPU) is critical when running non-trivial user programs.
    \item \textbf{Toolchain Limitations}: Educational OS environments often lack "standard" features (like \texttt{\%f} in printf or \texttt{mmap}), requiring creative workarounds.
    \item \textbf{Debugging}: Analyzing raw trap causes (\texttt{scause}, \texttt{sepc}) is essential for diagnosing low-level crashes like illegal instructions or unhandled page faults.
\end{itemize}

\end{document}
